\chapter{Euphoria}

\section*{Definition}
Euphoria is an exaggerated, intense feeling of well-being, elation, or happiness that is disproportionate to the situation. Pathological euphoria is sustained, contextually inappropriate, or occurs without a clear positive trigger, often indicating an underlying medical or psychiatric condition.

\section*{What Happens in Your Body}
Euphoria is driven by excessive dopaminergic and opioidergic signaling in the mesolimbic reward pathway (nucleus accumbens, ventral tegmental area, prefrontal cortex). Reduced prefrontal inhibitory control allows unchecked activation of these circuits, while serotonergic and GABAergic deficits fail to modulate the emotional response.

\section*{Causes}
\begin{itemize}
  \item \textbf{Psychiatric:} Bipolar I disorder (manic episode), bipolar II disorder (hypomanic episode), cyclothymia
  \item \textbf{Substance-induced:} Stimulants (cocaine, amphetamines), opioids, cannabis, alcohol, MDMA, hallucinogens, anabolic steroids
  \item \textbf{Neurological:} Multiple sclerosis (especially with frontal lesions), stroke, traumatic brain injury (orbitofrontal), brain tumors, epilepsy (simple partial seizures), narcolepsy
  \item \textbf{Endocrine/metabolic:} Hyperthyroidism, Cushing syndrome, adrenal tumors (pheochromocytoma)
  \item \textbf{Medication-induced:} Corticosteroids, levodopa, bupropion, thyroid hormone replacement, methylphenidate
\end{itemize}

\section*{When to See a Doctor}
See your doctor if:
\begin{itemize}
  \item Euphoria is uncharacteristic, sustained, and accompanied by grandiosity, decreased need for sleep, or reckless behavior
  \item There is rapid onset with confusion, head injury, or focal neurological signs
  \item The individual poses a risk due to poor judgment or impulsive decisions
\end{itemize}

\section*{Related Symptoms}
\begin{itemize}
  \item Mania, hypomania, irritability, pressured speech, flight of ideas, grandiosity, decreased need for sleep, hyperactivity, hypersexuality, substance use
\end{itemize}
\textbf{Important:} Euphoria in the context of bipolar disorder is often ego-syntonic, making the individual reluctant to seek treatment.

\section*{Self-Care / Home Management}
\begin{itemize}
  \item Avoid stimulants, alcohol, and other psychoactive substances
  \item Maintain a regular sleep schedule and track sleep duration
  \item Engage trusted contacts to provide perspective on behavior changes
  \item Use a mood chart to track mood, energy, and sleep patterns
\end{itemize}

\section*{How Your Doctor Will Evaluate You}
\begin{itemize}
  \item Clinical history assessing duration, context, associated symptoms, and functional impact
  \item Mental status examination: mood, affect, thought content (grandiosity), thought process (pressured, flight of ideas), judgment
  \item Collateral history from family or close contacts
  \item Drug toxicology screen, thyroid function tests, and neurological evaluation if indicated
  \item Distinguish mania from hypomania using DSM-5 criteria (duration, functional impairment)
\end{itemize}

\section*{Treatment Options}
\begin{itemize}
  \item First-line for bipolar mania: mood stabilizers (lithium, valproate, lamotrigine) and/or atypical antipsychotics (quetiapine, olanzapine, aripiprazole)
  \item Eliminate causative medications or substances
  \item Hospitalization if severe mania with psychosis, aggression, or risk-taking behavior
  \item Long-term maintenance with mood stabilizers to prevent recurrence
\end{itemize}

\section*{References}
\begin{thebibliography}{99}
\bibitem{euphoria:ref1} American Psychiatric Association. ``Diagnostic and Statistical Manual of Mental Disorders.'' 5th ed, Text Revision. APA, 2022.
\bibitem{euphoria:ref2} Goodwin FK, Jamison KR. ``Manic-Depressive Illness: Bipolar Disorders and Recurrent Depression.'' 2nd ed. Oxford University Press, 2007.
\bibitem{euphoria:ref3} Miklowitz DJ, Johnson SL. ``The Psychopathology and Treatment of Bipolar Disorder.'' Annu Rev Clin Psychol. 2006;2:199--235.
\bibitem{euphoria:ref4} Phillips ML, Kupfer DJ. ``Bipolar Disorder Diagnosis: Challenges and Future Directions.'' Lancet. 2013;381(9878):1663--1671.
\bibitem{euphoria:ref5} Coryell W, Fiedorowicz JG, Solomon DA, et al. ``Euphoric and Dysphoric Mania: A 20-Year Prospective Study.'' J Clin Psychiatry. 2018;79(1):17m11527.
\end{thebibliography}
